% Escreva aqui a solução da questão 9.

\textcolor{red}{ VERSÃO 1 } %Rafaela

Sejam $a, b, c, d \in \R_+^\ast$. Aplicando a desigualdade das médias
aritmética e geométrica nos termos $a , b, c ,d$, temos:
\begin{align*}
    \sqrt[4]{a\cdot b\cdot c\cdot d} \leq \dfrac{a+b+c+d}{4}.\tag{I} \\
\end{align*}

Agora, aplicando a desigualdade das médias harmônica e geométrica nos termos
$a, b, c, d$, temos:
\begin{align*}
    \dfrac{4}{\dfrac{1}{a}+\dfrac{1}{b}+\dfrac{1}{c}+\dfrac{1}{d}} \leq \sqrt[4]{a\cdot b\cdot c\cdot d}. \tag{II}
\end{align*}

Por transitividade de (I) e (II), temos:
\begin{align*}
    %            \dfrac{4}{\dfrac{1}{a}+\dfrac{1}{b}+\dfrac{1}{c}+\dfrac{1}{d}} \leq \sqrt[4]{a\cdot b\cdot c\cdot d} \leq  \dfrac{a+b+c+d}{4} \\
    \dfrac{4}{\dfrac{1}{a}+\dfrac{1}{b}+\dfrac{1}{c}+\dfrac{1}{d}}           & \leq \dfrac{a+b+c+d}{4}                                                                                                                                      \\
    \implies \dfrac{16}{\dfrac{1}{a}+\dfrac{1}{b}+\dfrac{1}{c}+\dfrac{1}{d}} & \leq a+b+c+d                                                                                                                                                 \\
    \implies 16                                                              & \leq (a+b+c+d) \cdot \left(\dfrac{1}{a}+\dfrac{1}{b}+\dfrac{1}{c}+\dfrac{1}{d}\right). & \text{$\left(\dfrac 1 a+\dfrac 1 b+\dfrac 1 c+\dfrac 1 d>0\right)$}
\end{align*}

Portanto, $(a+b+c+d) \cdot
    \left(\dfrac{1}{a}+\dfrac{1}{b}+\dfrac{1}{c}+\dfrac{1}{d}\right) \geq 16$, para
quaisquer $a, b, c, d \in \R_+^\ast$.

%\bigskip

\textcolor{red}{ VERSÃO 2} %Rafaela

Sejam $a, b, c, d \in \R_+^\ast$. Aplicando a desigualdade das médias
aritmética e geométrica nos termos $a , b, c ,d$, temos:
\begin{align*}
    \sqrt[4]{a\cdot b\cdot c\cdot d} \leq \dfrac{a+b+c+d}{4}.\tag{I} \\
\end{align*}

Agora, aplicando a desigualdade das médias harmônica e geométrica nos termos
$a, b, c, d$, temos:
\begin{align*}
    \dfrac{4}{\dfrac{1}{a}+\dfrac{1}{b}+\dfrac{1}{c}+\dfrac{1}{d}} \leq \sqrt[4]{a\cdot b\cdot c\cdot d} & \implies \dfrac{\dfrac{1}{a}+\dfrac{1}{b}+\dfrac{1}{c}+\dfrac{1}{d}}{4} \geq \dfrac{1}{\sqrt[4]{a\cdot b\cdot c\cdot d}}            \\
                                                                                                         & \implies \dfrac{1}{\sqrt[4]{a\cdot b\cdot c\cdot d}} \leq  \dfrac{\dfrac{1}{a}+\dfrac{1}{b}+\dfrac{1}{c}+\dfrac{1}{d}}{4}. \tag{II}
\end{align*}

``Multiplicando'' (I) e (II) (todos os termos são positivos), temos:
\begin{align*}
    \left(\sqrt[4]{a\cdot b\cdot c\cdot d}\right) \cdot \left(\dfrac{1}{\sqrt[4]{a\cdot b\cdot c\cdot d}}\right) & \leq \left(\dfrac{a+b+c+d}{4}\right) \cdot \left(\dfrac{\dfrac{1}{a}+\dfrac{1}{b}+\dfrac{1}{c}+\dfrac{1}{d}}{4}\right) \\
    \implies 1                                                                                                   & \leq \dfrac{(a+b+c+d)\cdot \left(\dfrac{1}{a}+\dfrac{1}{b}+\dfrac{1}{c}+\dfrac{1}{d}\right)}{16}                       \\
    %       \implies 1 \cdot 16 &\leq \dfrac{(a+b+c+d)\cdot \left(\dfrac{1}{a}+\dfrac{1}{b}+\dfrac{1}{c}+\dfrac{1}{d}\right)}{16} \cdot 16 \\
    \implies 16                                                                                                  & \leq (a+b+c+d) \cdot \left(\dfrac{1}{a}+\dfrac{1}{b}+\dfrac{1}{c}+\dfrac{1}{d}\right).
\end{align*}

Portanto, $(a+b+c+d) \cdot
    \left(\dfrac{1}{a}+\dfrac{1}{b}+\dfrac{1}{c}+\dfrac{1}{d}\right) \geq 16$, para
quaisquer $a, b, c, d \in \R_+^\ast$.

%\bigskip

%\textcolor{red}{ VERSÃO 3} % autor

%\bigskip

\bigskip

\textcolor{red}{CHAVE DE PONTUAÇÃO:
    \begin{itemize}
        \item  Obtenção das 2 desigualdades das médias \emph{(1,0 ponto)};
        \item Manipulação das inequações \emph{(1,0 ponto)};
        \item  Redação em si (desenho, notação e organização lógica) \emph{(1,0 ponto)}.
    \end{itemize} }